\subsection{Conclusion}

This thesis Discussed RAGPoison Lab made specfically to respond to the growing issue of recommender systems tampering when involving LLM and RAG. It comes from a security concern, How much impact this type of red-teaming can have on user in real life context on their decision planning. Instead of attacking a live service we recreated a local replica of such an environnement with a local dataset and created a modular workflow that can study every step of the workflow through multiple attack families, policies and more.

We benchmarked different models at various intensities ranking methods and got verying results but that all still lead to one same concern. This type of attack is real and can impact unknowing users in more cases than one. Such vulerability should not be limited to benign setups like this one where the only impact are the movies watched but rather should be looked at from the perspective of other cases like consummer fields where a user can be convinced to buy a product that would go against his actual interests.

The final matrix experiment shows that the specific attacks will lead to different interpretations. Targeted ppromotion is best undertood when looking at the item that is targeted to enter top-$k$ results, prompt injection also has a similar path but utilises the models weakness for jailbreaking than the other attacks. untargeted degredation should always be looked at through the quality metrics. we also understand that these systems dont seem to be ready for consummer use as all these attacks still get tothe same conclusion: the attack suceeded greatly for most users. 

This thesis still doesn't claim that its a representation of all similar systems and that it proves a universal issue that come LLM-powered recommender system. It also assumes that at this scale the results might not be comparablle to ones of massive systems a consummer might encounter on the internet. It is a small scale local based and built from scratch small simulation tool for RAG recommenders that is meant to be safe and reproducible.

\subsection{Future Work}

Future attempts should use a different Dataset, with larger and richer data like new or conversational recommendation. second, it should utilise all the tools that are ready to use in RAGPoison, which will take more than a realistic few months of continuous runs to achieve but will allow to have a wider picture of the subject. Third, Future work should attempt defense mechanism against such attack and try to achieve a real difference accross all attack families.

Adding more tools, more attack type could also improve the project, extend the abilities of python SDK to allow more detailed dashboard and more in the objective of broadening the scope of this benchmark tool to be acome a real security first tool for the study of attacks and solution generaton